A \emph{(distinct/incongruent) covering system} of $\mathbb{Z}$ is a finite collection of congruences
\[
x \equiv a_i \pmod{m_i}\qquad (1\le i\le k),
\]
with integers $m_i\ge 2$ that are \emph{pairwise distinct}, such that every $x\in\mathbb{Z}$ satisfies at least one of the congruences.

The literal wording in the problem statement does not explicitly say ``distinct moduli''.  However, if repetition is allowed then the problem becomes trivial: since $4 = 5-1$, the four congruences $x\equiv 0,1,2,3\pmod 4$ already cover $\mathbb{Z}$ using moduli of the form $p-1$ with $p\ge 5$.  Thus the only nontrivial interpretation is the distinct-modulus version.

\medskip
\noindent\textbf{Corrected statement (standard in covering-systems theory).}
\emph{Does there exist a distinct covering system of $\mathbb{Z}$ in which every modulus is of the form $p-1$ for some prime $p\ge 5$?}
